\textbf{Время:} резерв

\textbf{Тезис:} текущая трёхмерная корреляция доступна также в выбранном интервале быстроты пары.

\textbf{На экране:} проекции out, side, long; общие условия и интервал y внизу.

\textbf{Текст:} Это другой отбор текущих модельных пар: вместо узкого интервала k-T выбран интервал быстроты от нуля до 0,2. Заряд отрицательный, центральность 0–10 процентов. Быстрота задана в системе центра масс столкновения; радиусы определяются в LCMS каждой пары.

Чёрные маркеры и красные линии читаются так же, как на основном слайде 26. Это отдельная выборка, поэтому её параметры не нужно смешивать с параметрами узкого k-T-интервала.

\textbf{Переход:} Повторив отбор в соседних интервалах быстроты, получаем зависимости радиусов от y.

\textbf{Справка:} cfs\_neg\_0-10\_0.0\_0.2.jpg. Подписи файла подтверждают neg, 0–10\%, y=[0.0,0.2]. Напечатаны \(R_{\rm out}=6.051\pm0.008\) фм, \(R_{\rm side}=4.616\pm0.005\) фм, \(R_{\rm long}=4.294\pm0.005\) фм, \(\lambda=0.889\pm0.001\), \(\chi^2/\mathrm{ndf}=1.214\). Это прочитанные аннотации аппроксимации; новые погрешности не рассчитывались. Полный kT-отбор и проекционное окно в JPEG не заданы. Нуль Rout–long не трактуется как самостоятельный физический результат.

\textbf{Источники:} предоставленный автором JPG, подтверждения LCMS и системы центра масс; текущие манифесты.
